% GENERATED FILE. Edit canonical lessons, then run npm run generate:books.
% canonical-source-hash: fnv1a64:18da5e82f2bd6fed
% canonical-lessons: ML-C426-bahiskarikkuka, ML-C426-avasyappetuka, ML-C426-nivedanam, ML-C426-samarppikkuka, ML-C426-harji

\chapter{Demands}
\label{ch:ml-demands}
\hlchaptermodality{426}

\begin{chapteropening}
\textbf{By the end of this chapter:} \emph{I can say to boycott, to demand, a memorandum, to submit and a petition.}

Complete the last lesson of chapter 426: I can say to boycott, to demand, a memorandum, to submit and a petition.
\end{chapteropening}

\section[\texorpdfstring{\ml{ബഹിഷ്കരിക്കുക} (bahiṣkarikkuka)}{ബഹിഷ്കരിക്കുക (bahiṣkarikkuka)}]{\ml{ബഹിഷ്കരിക്കുക} (bahiṣkarikkuka) — to boycott}

\label{lesson:ML-C426-bahiskarikkuka}



\begin{quote}
Before the new one: say the Malayalam for a general strike, then the Malayalam for a march.
\end{quote}

\subsection*{You'll want to know: \ml{ബഹിഷ്കരിക്കുക}}
\textbf{\ml{ബഹിഷ്കരിക്കുക}} — \emph{bahiṣkarikkuka} — \textquotedblleft{}to boycott\textquotedblright{}.

\begin{quote}
\ml{വിദ്യാർത്ഥികൾ} \ml{ക്ലാസുകൾ} \ml{ബഹിഷ്കരിച്ചു}.
\end{quote}

\begin{quote}
\emph{vidyārtthikaḷ klāsukaḷ bahiṣkariccu.}
\end{quote}

\begin{quote}
The students boycotted classes.
\end{quote}

Past: \ml{ബഹിഷ്കരിച്ചു}, \emph{bahiṣkariccu}.

The noun is \ml{ബഹിഷ്കരണം}, a boycott.

\begin{cousinweb}
\textbf{\ml{ബഹിഷ്കരിക്കുക}} (\emph{bahiṣkarikkuka}) is built from Sanskrit \emph{bahiṣkāra}, \textquotedblleft{}expulsion, removal\textquotedblright{} (Monier-Williams), from \emph{bahis}, \textquotedblleft{}outside\textquotedblright{}, and \emph{kṛ}, \textquotedblleft{}to make\textquotedblright{}, with the Malayalam verb ending: to put a thing outside.
\end{cousinweb}

\begin{grammarlens}[title={What you've built}]
One more word for giving reasons and opinions.
\end{grammarlens}

\subsection*{Your turn}
\begin{itemize}
\raggedright
  \item \emph{Say it:} \emph{bahiṣkarikkuka}
  \item \emph{Say it:} \emph{bahiṣkarikkuka}, once more
  \item \emph{Say it:} \emph{vidyārtthikaḷ klāsukaḷ bahiṣkariccu.} — The students boycotted classes.
  \item \emph{Recall:} say the Malayalam for a general strike, then the Malayalam for a march, then say \emph{bahiṣkarikkuka} again
\end{itemize}

\begin{tcolorbox}[breakable,colback=teal!4,colframe=teal!35!black,title={Before you move on}]
What does \ml{ബഹിഷ്കരിക്കുക} mean? (\textquotedblleft{}To boycott\textquotedblright{}.) Say it once more.
\end{tcolorbox}
\section[\texorpdfstring{\ml{ആവശ്യപ്പെടുക} (āvaśyappeṭuka)}{ആവശ്യപ്പെടുക (āvaśyappeṭuka)}]{\ml{ആവശ്യപ്പെടുക} (āvaśyappeṭuka) — to demand; to request}

\label{lesson:ML-C426-avasyappetuka}



\begin{quote}
Before the new one: say the Malayalam for a march, then the Malayalam for to boycott.
\end{quote}

\subsection*{You'll want to know: \ml{ആവശ്യപ്പെടുക}}
\textbf{\ml{ആവശ്യപ്പെടുക}} — \emph{āvaśyappeṭuka} — \textquotedblleft{}to demand; to request\textquotedblright{}.

\begin{quote}
\ml{തൊഴിലാളികൾ} \ml{കൂടുതൽ} \ml{ശമ്പളം} \ml{ആവശ്യപ്പെട്ടു}.
\end{quote}

\begin{quote}
\emph{toḻilāḷikaḷ kūṭutal śampaḷaṁ āvaśyappeṭṭu.}
\end{quote}

\begin{quote}
The workers demanded more pay.
\end{quote}

Past: \ml{ആവശ്യപ്പെട്ടു}, \emph{āvaśyappeṭṭu}.

Built on \ml{ആവശ്യം}, a need, with -\ml{പ്പെടുക}.

\begin{cousinweb}
\textbf{\ml{ആവശ്യപ്പെടുക}} (\emph{āvaśyappeṭuka}) is built on \ml{ആവശ്യം}, \textquotedblleft{}a need\textquotedblright{}, from Sanskrit \emph{āvaśyaka}, \textquotedblleft{}necessary\textquotedblright{} (Monier-Williams), with -\ml{പ്പെടുക} from native \ml{പെടുക} (DEDR 3853): to declare a thing needed.
\end{cousinweb}

\begin{grammarlens}[title={What you've built}]
One more word for giving reasons and opinions.
\end{grammarlens}

\subsection*{Your turn}
\begin{itemize}
\raggedright
  \item \emph{Say it:} \emph{āvaśyappeṭuka}
  \item \emph{Say it:} \emph{āvaśyappeṭuka}, once more
  \item \emph{Say it:} \emph{toḻilāḷikaḷ kūṭutal śampaḷaṁ āvaśyappeṭṭu.} — The workers demanded more pay.
  \item \emph{Recall:} say the Malayalam for a march, then the Malayalam for to boycott, then say \emph{āvaśyappeṭuka} again
\end{itemize}

\begin{tcolorbox}[breakable,colback=teal!4,colframe=teal!35!black,title={Before you move on}]
What does \ml{ആവശ്യപ്പെടുക} mean? (\textquotedblleft{}To demand; to request\textquotedblright{}.) Say it once more.
\end{tcolorbox}
\section[\texorpdfstring{\ml{നിവേദനം} (nivēdanaṁ)}{നിവേദനം (nivēdanaṁ)}]{\ml{നിവേദനം} (nivēdanaṁ) — a memorandum; a petition}

\label{lesson:ML-C426-nivedanam}



\begin{quote}
Before the new one: say the Malayalam for to boycott, then the Malayalam for to demand.
\end{quote}

\subsection*{You'll want to know: \ml{നിവേദനം}}
\textbf{\ml{നിവേദനം}} — \emph{nivēdanaṁ} — \textquotedblleft{}a memorandum; a petition\textquotedblright{}.

\begin{quote}
\ml{നാട്ടുകാർ} \ml{മന്ത്രിക്ക്} \ml{ഒരു} \ml{നിവേദനം} \ml{നൽകി}.
\end{quote}

\begin{quote}
\emph{nāṭṭukār mantrikkŭ oru nivēdanaṁ nalki.}
\end{quote}

\begin{quote}
The local people gave the minister a memorandum.
\end{quote}

It is given, \ml{നൽകുക}, to an official.

\begin{cousinweb}
\textbf{\ml{നിവേദനം}} (\emph{nivēdanaṁ}) is borrowed from Sanskrit \emph{nivedana}, \textquotedblleft{}announcing, making known\textquotedblright{} (Monier-Williams): a memorandum makes a need known.
\end{cousinweb}

\begin{grammarlens}[title={What you've built}]
One more word for giving reasons and opinions.
\end{grammarlens}

\subsection*{Your turn}
\begin{itemize}
\raggedright
  \item \emph{Say it:} \emph{nivēdanaṁ}
  \item \emph{Say it:} \emph{nivēdanaṁ}, once more
  \item \emph{Say it:} \emph{nāṭṭukār mantrikkŭ oru nivēdanaṁ nalki.} — The local people gave the minister a memorandum.
  \item \emph{Recall:} say the Malayalam for to boycott, then the Malayalam for to demand, then say \emph{nivēdanaṁ} again
\end{itemize}

\begin{tcolorbox}[breakable,colback=teal!4,colframe=teal!35!black,title={Before you move on}]
What does \ml{നിവേദനം} mean? (\textquotedblleft{}A memorandum; a petition\textquotedblright{}.) Say it once more.
\end{tcolorbox}
\section[\texorpdfstring{\ml{സമർപ്പിക്കുക} (samarppikkuka)}{സമർപ്പിക്കുക (samarppikkuka)}]{\ml{സമർപ്പിക്കുക} (samarppikkuka) — to submit}

\label{lesson:ML-C426-samarppikkuka}



\begin{quote}
Before the new one: say the Malayalam for to demand, then the Malayalam for a memorandum.
\end{quote}

\subsection*{You'll want to know: \ml{സമർപ്പിക്കുക}}
\textbf{\ml{സമർപ്പിക്കുക}} — \emph{samarppikkuka} — \textquotedblleft{}to submit\textquotedblright{}.

\begin{quote}
\ml{അവൾ} \ml{അപേക്ഷ} \ml{ഇന്നലെ} \ml{സമർപ്പിച്ചു}.
\end{quote}

\begin{quote}
\emph{avaḷ apēkṣa innale samarppiccu.}
\end{quote}

\begin{quote}
She submitted the application yesterday.
\end{quote}

Past: \ml{സമർപ്പിച്ചു}, \emph{samarppiccu}.

Also to dedicate: \ml{പുസ്തകം} \ml{അമ്മയ്ക്ക്} \ml{സമർപ്പിച്ചു}, he dedicated the book to his mother.

\begin{cousinweb}
\textbf{\ml{സമർപ്പിക്കുക}} (\emph{samarppikkuka}) is built from Sanskrit \emph{samarpaṇa}, \textquotedblleft{}placing upon; delivering over\textquotedblright{} (CDIAL 13191), with the Malayalam verb ending -\ml{ിക്കുക}.
\end{cousinweb}

\begin{grammarlens}[title={What you've built}]
One more word for giving reasons and opinions.
\end{grammarlens}

\subsection*{Your turn}
\begin{itemize}
\raggedright
  \item \emph{Say it:} \emph{samarppikkuka}
  \item \emph{Say it:} \emph{samarppikkuka}, once more
  \item \emph{Say it:} \emph{avaḷ apēkṣa innale samarppiccu.} — She submitted the application yesterday.
  \item \emph{Recall:} say the Malayalam for to demand, then the Malayalam for a memorandum, then say \emph{samarppikkuka} again
\end{itemize}

\begin{tcolorbox}[breakable,colback=teal!4,colframe=teal!35!black,title={Before you move on}]
What does \ml{സമർപ്പിക്കുക} mean? (\textquotedblleft{}To submit\textquotedblright{}.) Say it once more.
\end{tcolorbox}
\section[\texorpdfstring{\ml{ഹർജി} (harji)}{ഹർജി (harji)}]{\ml{ഹർജി} (harji) — a petition (to a court)}

\label{lesson:ML-C426-harji}



\begin{quote}
Before the new one: say the Malayalam for a memorandum, then the Malayalam for to submit.
\end{quote}

\subsection*{You'll want to know: \ml{ഹർജി}}
\textbf{\ml{ഹർജി}} — \emph{harji} — \textquotedblleft{}a petition (to a court)\textquotedblright{}.

\begin{quote}
\ml{അവർ} \ml{കോടതിയിൽ} \ml{ഒരു} \ml{ഹർജി} \ml{നൽകി}.
\end{quote}

\begin{quote}
\emph{avar kōṭatiyil oru harji nalki.}
\end{quote}

\begin{quote}
They filed a petition in court.
\end{quote}

Plural: \ml{ഹർജികൾ}.

\begin{cousinweb}
\textbf{\ml{ഹർജി}} (\emph{harji}) is usually traced to Hindustani \emph{arzī}, \textquotedblleft{}a petition\textquotedblright{}, which came from Arabic through Persian; where the Malayalam initial \ml{ഹ}- comes from is not settled.
\end{cousinweb}

\begin{grammarlens}[title={What you've built}]
One more word for giving reasons and opinions.
\end{grammarlens}

\subsection*{Your turn}
\begin{itemize}
\raggedright
  \item \emph{Say it:} \emph{harji}
  \item \emph{Say it:} \emph{harji}, once more
  \item \emph{Say it:} \emph{avar kōṭatiyil oru harji nalki.} — They filed a petition in court.
  \item \emph{Recall:} say the Malayalam for a memorandum, then the Malayalam for to submit, then say \emph{harji} again
\end{itemize}

\begin{tcolorbox}[breakable,colback=teal!4,colframe=teal!35!black,title={Before you move on}]
What does \ml{ഹർജി} mean? (\textquotedblleft{}A petition (to a court)\textquotedblright{}.) Say it once more.
\end{tcolorbox}
